% !TeX root = ../../tesis.tex

\section{Servidor Proxy de Geometrías para Tableau (\termino{GeoProxy})}
\label{anx:tg-geoproxy}

El conector \termino{WDC} transmite filas tabulares pero no puede manejar
geometrías complejas como \texttt{MultiLineString} o \texttt{MultiPolygon}. Para
esas capas, Tableau Desktop dispone de un conector de \termino{Spatial File} que
apunta a un archivo \geojson{} accesible por \textsc{http}. El servidor
\textbf{GeoProxy} (\texttt{geo\_proxy/main.py}) opera como puente: re-expone los
endpoints geoespaciales de \textbf{Apimetro} y \textbf{VFTModel} como
\textsc{url}s locales que Tableau reconoce como fuentes espaciales sin necesidad
de guardar archivos en disco.

El servidor corre sobre \textbf{FastAPI} en el puerto~\texttt{5050} y ofrece
seis rutas organizadas en dos grupos. Su único componente de lógica es la
función \texttt{\_proxy}, que realiza un \textsc{get} asíncrono al servicio
externo y reenvía el cuerpo con tipo \textsc{mime} \texttt{application/geo+json}
, preservando el código de estado original. El Cuadro~
\ref{tab:tg-geoproxy-endpoints} lista todos los endpoints disponibles.

\begin{table}[htbp]
    \centering
    \caption[Endpoints del \termino{GeoProxy} para Tableau Spatial File]{Rutas
    del servidor \textbf{GeoProxy} disponibles como fuentes espaciales
    en Tableau Desktop mediante el conector \termino{Spatial File}.}
    \label{tab:tg-geoproxy-endpoints}
    \begin{tabularx}{\textwidth}{|l|l|X|}
        \hline
        \textbf{URL local (:5050)} & \textbf{Geometría} & \textbf{Endpoint origen} \\
        \hline
        \texttt{/apimetro/lineas.geojson}     & MultiLineString & \texttt{Apimetro/movilidad/mapas/geojsonLinea} \\
        \texttt{/apimetro/poligonos.geojson}  & MultiPolygon    & \texttt{Apimetro/movilidad/mapas/geojsonPoligono} \\
        \texttt{/apimetro/estaciones.geojson} & Point           & \texttt{Apimetro/movilidad/mapas/geojsonEstacion} \\
        \texttt{/vftmodel/cobertura.geojson}  & Polygon         & \texttt{VFTModel/api/v1/network/geolayers/coverage} \\
        \texttt{/vftmodel/capillary.geojson}  & Point           & \texttt{VFTModel/api/v1/network/geolayers/capillary} \\
        \texttt{/vftmodel/detour.geojson}     & Point           & \texttt{VFTModel/api/v1/network/geolayers/detour} \\
        \hline
    \end{tabularx}
    \fuentefigura{Elaboración propia con base en \texttt{geo\_proxy/main.py}
    (Transport-gis-zmvm-mjg, 2026).}
\end{table}

La configuración del servidor se gestiona mediante dos variables de entorno
definidas en \texttt{geo\_proxy/.env} (plantilla en \texttt{.env.example}):
\texttt{APIMETRO\_URL} (por defecto \texttt{http://localhost:8080}) y
\texttt{VFTMODEL\_URL} (por defecto \texttt{http://localhost:8000}). La
secuencia de arranque desde la raíz del repositorio es:

\begin{lstlisting}[style=estiloBase, language=bash,
  caption={Arranque del servidor \textbf{GeoProxy} en el puerto~5050.},
  label=lst:tg-geoproxy-arranque,
  basicstyle=\small\ttfamily]
cd geo_proxy
cp .env.example .env          # ajustar APIMETRO_URL y VFTMODEL_URL
pip install -r requirements.txt
uvicorn main:app --port 5050 --reload
\end{lstlisting}

Una vez activo, Tableau Desktop puede conectarse mediante
\textbf{Conectar~$\to$ Más\ldots~$\to$ Archivo espacial} escribiendo la
\textsc{url} del endpoint deseado (por ejemplo,
\texttt{http://localhost:5050/apimetro/lineas.geojson}).

El Código~\ref{lst:tg-geoproxy-proxy} presenta la función \texttt{\_proxy} y el
Código~\ref{lst:tg-geoproxy-endpoint} un endpoint representativo.

\lstinputlisting[ style=estiloPython, firstline=57, lastline=65, caption={Función auxiliar \texttt{\_proxy}: proxy asíncrono que reenvía el cuerpo con tipo \texttt{application/geo+json}.}, label=lst:tg-geoproxy-proxy, breaklines=true, basicstyle=\tiny\ttfamily ]{Anexos/Transport-gis-mjg/geo_proxy/main.txt}

\lstinputlisting[ style=estiloPython, firstline=68, lastline=79, caption={Endpoint \texttt{/apimetro/lineas.geojson}: ejemplo representativo de los seis endpoints del servidor.}, label=lst:tg-geoproxy-endpoint, breaklines=true, basicstyle=\tiny\ttfamily ]{Anexos/Transport-gis-mjg/geo_proxy/main.txt}
